\section*{Governing equations}
\label{sec:strong-forms-nd}

\paragraph{Smooth regularization operators.}
To ensure numerical stability and differentiability, we employ a consistent set of smooth regularization operators parameterized by a smoothing width $\varepsilon > 0$. We define the smooth absolute value as $|x|_\varepsilon := \sqrt{x^2+\varepsilon^2}-\varepsilon$. Based on this, we define smooth maximum, minimum, and clamping operators:
\begin{align} \label{eq:smooth-max-nd}
\operatorname{smax}(x,y;\varepsilon) &:= \tfrac{1}{2}(x+y+|x-y|_\varepsilon), \nonumber \\
\operatorname{smin}(x,y;\varepsilon) &:= \tfrac{1}{2}(x+y-|x-y|_\varepsilon), \\
\operatorname{sclamp}(x,a,b;\varepsilon) &:= \operatorname{smin}(\operatorname{smax}(x,a;\varepsilon),b;\varepsilon). \nonumber
\end{align}

\paragraph{Hard domain guards (numerical clipping).}
Some intermediate expressions used in purely numerical subroutines (notably closed-form eigensolvers and invariant mappings involving $\arccos(\cdot)$ or $\ln(\cdot)$) require \emph{strict} range preservation to avoid NaNs. For these numerical guards---which are not part of the constitutive or biological model assumptions---we use the sharp pointwise clipping operator
\[
\operatorname{clip}(x,a,b):=\min\{\max\{x,a\},b\},
\]
implemented as sharp pointwise min/max operations. We intentionally keep these guards sharp to guarantee domain safety.

\paragraph{Units and scaling.}
We use the unit system: length [mm], time [day], stress/energy [MPa], and density [g/cm$^3$]. The stimulus $S$ is dimensionless, the apparent density $\rho$ is in g/cm$^3$, and fabric quantities ($\mathbf{A}$, $\mathbf{L}$) are dimensionless. The boundary partition is $\partial\Omega=\Gamma_D\cup\Gamma_N$ with outward unit normal $\mathbf{n}$.

\subsection*{Mechanical submodel}
\label{subsec:mech-nd}
\paragraph{Kinematics and equilibrium.}
We adopt small-strain kinematics with the symmetric gradient
\begin{equation}\label{eq:mech-eps-nd}
\boldsymbol{\varepsilon}(\mathbf{u}) := \tfrac12\big(\nabla\mathbf{u}+\nabla\mathbf{u}^\top\big)\in\mathbb{R}^{d\times d}_{\rm sym},\qquad d\in\{2,3\}.
\end{equation}
Under quasi-static conditions the balance of linear momentum reads
\begin{equation}\label{eq:mech-strong-nd}
-\nabla\!\cdot\!\boldsymbol{\sigma}(\mathbf{u},\rho,\mathbf{L})=\mathbf{0}\ \ \text{in }\Omega,\qquad
\mathbf{u}=\mathbf{0}\ \ \text{on }\Gamma_D,\qquad
\boldsymbol{\sigma}(\mathbf{u},\rho,\mathbf{L})\,\mathbf{n}=\mathbf{t}\ \ \text{on }\Gamma_N.
\end{equation}

\paragraph{Constitutive structure.}
We adopt the Zysset-Curnier orthotropic material model \cite{Zysset1995_AlternativeModelAnisotropic}, which defines the stiffness tensor $\mathbb{C}$ based on the apparent density $\rho$ and the fabric eigenvalue ratios $m_1, m_2, m_3$ obtained from the log-fabric field $\mathbf{L}$. The fabric tensor is recovered from $\mathbf{L}$ via the matrix exponential $\mathbf{A} = \exp(\mathbf{L})$.
\rev{The law is used here as an apparent-level constitutive closure: unlike bottom-up osteon-scale elasto-damage models \revcite{Gizzi2022_OsteonMechanics}, it does not resolve explicit microstructural damage mechanisms, but instead encodes anisotropy through the evolving fabric field.}
The orthotropic engineering constants are given by:
\begin{equation}\label{eq:zysset-curnier}
\begin{aligned}
    E_i &= E_0 \left(\frac{\widetilde{\rho}}{\rho_{\mathrm{ref}}}\right)^{k(\widetilde{\rho})} m_i^{p_E}, \\
    G_{ij} &= \frac{E_0}{2(1+\nu_0)} \left(\frac{\widetilde{\rho}}{\rho_{\mathrm{ref}}}\right)^{k(\widetilde{\rho})} (m_i m_j)^{p_G/2}, \\
    \nu_{ij} &= \nu_0 \sqrt{\frac{E_i}{E_j}} \quad \text{for } i\neq j,
\end{aligned}
\end{equation}
where $E_0$ is a reference Young's modulus, $\nu_0$ is a (density-independent) Poisson ratio, $\rho_{\mathrm{ref}}$ is a reference density, and $p_E$ and $p_G$ are stiffness exponents controlling axial and shear scaling with fabric. The density-dependent exponent $k(\widetilde{\rho})$ transitions smoothly from a trabecular exponent $n_{\mathrm{trab}}$ to a cortical exponent $n_{\mathrm{cort}}$:
\begin{equation}\label{eq:variable-exponent}
k(\widetilde{\rho}) = n_{\mathrm{trab}} (1 - w(\widetilde{\rho})) + n_{\mathrm{cort}} w(\widetilde{\rho}),
\end{equation}
where $w(\widetilde{\rho})$ is a smooth interpolation function (smoothstep) transitioning from 0 to 1 over the range $[\rho_{\mathrm{trab}}^{\max}, \rho_{\mathrm{cort}}^{\min}]$. The density $\widetilde{\rho}$ is smoothly clamped to a minimum value $\rho_{\min}$ to ensure positive definiteness.
\paragraph{Orthotropic tensor construction and physical admissibility.}
The coefficients \eqref{eq:zysset-curnier} define engineering constants in the \emph{principal fabric frame}.
Let $\widehat{\mathbf{A}}$ denote the \emph{deviatoric (isotropic-scaling-removed) normalization} of the fabric tensor, defined by \eqref{eq:Ahat-blend-nd}. We denote by $\{\mathbf{a}_i\}_{i=1}^3$ the orthonormal eigenvectors of $\widehat{\mathbf{A}}$ and by $\{m_i\}_{i=1}^3$ its eigenvalues, which represent bounded \emph{anisotropy factors} with isotropic scaling removed in log space.
The stiffness operator entering $\boldsymbol{\sigma}=\mathbb{C}:\boldsymbol{\varepsilon}$ is conveniently expressed in projector form using the spectral decomposition
\(
\widehat{\mathbf{A}}=\sum_{i=1}^3 m_i\,\mathbf{P}_i
\)
with $\mathbf{P}_i:=\mathbf{a}_i\otimes\mathbf{a}_i$ in the fully anisotropic case. Conceptually, one may assemble $\mathbb{C}$ by forming the orthotropic compliance matrix $\mathbb{S}$ from $\{E_i,G_{ij},\nu_{ij}\}$ (Voigt notation), enforcing reciprocity ($\nu_{ji}E_i=\nu_{ij}E_j$), inverting $\mathbb{S}$, and rotating to the spatial frame.
This construction ensures that the resulting constitutive operator is \emph{conservative} (minor and major symmetries) and admits a standard strain energy density $U=\tfrac12\boldsymbol{\sigma}:\boldsymbol{\varepsilon}$.
We ensure \emph{strong ellipticity} through admissible parameter ranges and bounded anisotropy (via eigenvalue-ratio bounds), and the law reduces smoothly to the isotropic baseline as $\widehat{\mathbf{A}}\to\mathbf{I}$.

For later use (e.g.\ fabric-guided constitutive evaluation) we remove isotropic scaling in log space and bound the resulting eigenvalue ratios. Let $\mathrm{sym}(\mathbf{L})=\sum_{i=1}^3 \ell_i\,\mathbf{P}_i$ with $\bar\ell:=\tfrac13\sum_i \ell_i$ \rev{(where $\{\ell_i,\mathbf{P}_i\}$ denote the eigenvalues and spectral projectors obtained from the eigendecomposition of $\mathrm{sym}(\mathbf{L})$)} and define
\begin{equation}\label{eq:Ahat-blend-nd}
\widehat{\mathbf{A}} \;:=\; \sum_{i=1}^3 m_i\,\mathbf{P}_i,\qquad
m_i \;:=\; \exp\!\Big(\operatorname{sclamp}(\ell_i-\bar\ell, -d_{\max}, d_{\max}; \varepsilon)\Big),\qquad
d_{\max}:=\ln\!\Big(\max\{m_{\max},1/m_{\min}\}\Big),
\end{equation}
which is symmetric positive definite. If the bound is inactive then $\det(\widehat{\mathbf{A}})=1$, while clamping enforces $m_i\in[\exp(-d_{\max}),\exp(d_{\max})]$. Robust evaluation of the spectral projectors near eigenvalue degeneracy and avoiding overflow in $\exp(\mathbf{L})$ require additional numerical safeguards; we use a stabilized Sylvester-projector construction and the above deviatoric-log clamping, see Appendix~\ref{sec:projector-robustness}.

\paragraph{Log-Euclidean Fabric Framework.}
To ensure the fabric tensor $\mathbf{A}$ remains symmetric and positive-definite (SPD) throughout the simulation, we evolve its matrix logarithm $\mathbf{L} = \ln \mathbf{A}$. The fabric tensor is recovered via $\mathbf{A} = \exp(\mathbf{L})$. This parameterization guarantees that $\mathbf{A}$ is always SPD for any real symmetric tensor $\mathbf{L}$, avoiding the need for ad-hoc projections or constraints, see the details in Appendix~\ref{sec:structure-preserving}.

\paragraph{Quantitative analysis and biological interpretation.}
Figure \ref{fig:constitutive_law} illustrates the behavior of the density-elasticity law. Panel (a) shows the stiffness $E(\rho)$ transitioning from a quadratic trabecular regime ($n_{\mathrm{trab}}=2.0$) to a near-linear cortical regime ($n_{\mathrm{cort}}=1.3$). This dual-power law captures the stiffening efficiency of dense bone while maintaining the rapid compliance increase in porous bone. Panel (b) explicitly plots the exponent $n(\rho)$, which varies smoothly between the two regimes via the cubic smoothstep function shown in (f). Panels (c) and (d) demonstrate the sensitivity of the stiffness curve to the transition thresholds $\rho_{\mathrm{trab}}^{\max}$ and $\rho_{\mathrm{cort}}^{\min}$; shifting these parameters allows for calibration to specific anatomical sites where the trabecular-cortical boundary may differ. Panel (e) highlights the effect of the exponents themselves; higher $n_{\mathrm{trab}}$ suppresses stiffness in low-density regions, effectively modeling the loss of connectivity in osteoporotic bone. The smooth lower clamp guarantees strictly positive stiffness even in highly porous regions, thus preventing mechanical degeneracy and preserving differentiability. The normalization \eqref{eq:Ahat-blend-nd} removes isotropic scaling in log space and bounds anisotropy through eigenvalue-ratio limits; when clamping is inactive, the normalization reduces to a unit-determinant descriptor.

\begin{figure}[htbp]
    \centering
    \safeincludegraphics[width=\textwidth]{constitutive_law.png}
    \caption{Constitutive Law Analysis. (a) Stiffness-density relationship $E(\rho)$ showing the transition from trabecular ($n=2.0$) to cortical ($n=1.3$) behavior. (b) The variable exponent $n(\rho)$. (c, d) Sensitivity of the transition to threshold parameters. (e) Effect of exponent choices on stiffness. (f) The cubic smoothstep function used for interpolation.}
    \label{fig:constitutive_law}
\end{figure}

\subsection*{Stimulus submodel}\label{subsec:stim-nd}
We represent the mechanotransductive drive by a scalar field $S:\Omega\times(0,T]\to\mathbb{R}$ that is \emph{generated} by mechanical activity and \emph{transported} through osteocytic/vascular pathways. \rev{By convention, positive values of $S$ correspond to net formation drive and negative values to net resorption drive; we note explicitly that this sign convention is a modelling assumption that directly maps the biological stimulus to the net remodeling outcome, thereby conflating the complex multi-cellular signalling cascade into a single scalar quantity.}

The stimulus obeys a diffusion--reaction balance with a smooth saturating (lazy-zone) mechanostat drive computed from a daily mechanical stimulus $\psi_{\mathrm{day}}$:
\begin{equation}\label{eq:stim-strong-nd}
\tau_S\,\partial_t S \;-\; \tau_S D_S\,\Delta S \;+\; S
\;=\; S_{\max}\,\tanh\!\Big(\delta_{\mathrm{eff}}(\psi_{\mathrm{day}},\rho)/\kappa_S\Big)
\quad\text{in }\Omega\times(0,T],
\end{equation}
supplemented with homogeneous Neumann boundary conditions and an initial state,
\begin{equation}\label{eq:stim-bc-ic-nd}
\nabla S\!\cdot\!\mathbf{n}=0\ \ \text{on }\partial\Omega\times(0,T],\qquad
S(\cdot,0)=S_0\ \ \text{in }\Omega.
\end{equation}
Here $\tau_S\ge 0$ is the stimulus time constant, $D_S\ge 0$ the stimulus diffusivity, $S_{\max}>0$ a cap on $|S|$, and $\kappa_S>0$ controls the saturation width in the $\tanh$ response. The drive is built from a normalized deviation in density-normalized strain energy,
\[
\hat{\psi}(x):=\frac{\psi_{\mathrm{day}}(x)}{\widetilde\rho(x)},\qquad
\hat{\psi}_{\mathrm{ref}}(\widetilde{\rho}):=\psi_{\mathrm{ref}}^{\mathrm{trab}}(1-w(\widetilde{\rho})) + \psi_{\mathrm{ref}}^{\mathrm{cort}}w(\widetilde{\rho}),\qquad
\delta(x):=\frac{\hat{\psi}(x)-\hat{\psi}_{\mathrm{ref}}(\widetilde{\rho})}{\hat{\psi}_{\mathrm{ref}}(\widetilde{\rho})}\,,
\]
where $\widetilde\rho$ is the smooth lower clamp \eqref{eq:smooth-max-nd} and $w(\widetilde{\rho})$ is the transition function \eqref{eq:variable-exponent}. To suppress spurious response to very small deviations we apply a smooth lazy-zone gate with threshold scale $\delta_0\ge 0$ (dimensionless):
\[
\delta_{\mathrm{eff}} \;:=\;
\begin{cases}
\delta\Big(1-\exp\!\big(-(\delta_{\mathrm{abs}}/\delta_0)^2\big)\Big), & \delta_0>0,\\
\delta, & \delta_0=0.
\end{cases}
\]
Here $\delta_{\mathrm{abs}}:=|\delta|_\varepsilon=\sqrt{\delta^2+\varepsilon^2}-\varepsilon$ is a smooth absolute value (dimensionless), consistent with the global smoothing width $\varepsilon$.
This choice acts as a smooth dead-zone-like gate: $\delta_{\mathrm{eff}}$ is strongly attenuated for $|\delta|\lesssim \delta_0$ and approaches $\delta$ for $|\delta|\gg \delta_0$.
The mechanical drive $\psi_{\mathrm{day}}$ is constructed as a cycle-weighted power mean of the local strain energy density (SED) across a discrete set of loading cases:
\begin{equation}\label{eq:stim-psi-nd}
\psi_{\mathrm{day}}(x) \;=\;
\left(\frac{\sum_i w_i\,\psi_i(x)^p}{\sum_i w_i}\right)^{1/p},\qquad
\psi_i(x):=\max\!\Big(0,\tfrac12\,\boldsymbol{\sigma}_i(x):\boldsymbol{\varepsilon}_i(x)\Big),
\end{equation}
where case weights $w_i$ are proportional to cycles-per-day (or relative daily frequency), and $p\ge 1$ biases the average toward peaks ($p=1$ gives a weighted mean, $p\to\infty$ approaches a max). The specific SED $\psi_{\mathrm{day}}/\widetilde\rho$ reflects sensing of strain energy per unit bone mass. The sharp max in $\psi_i$ enforces physical non-negativity of the strain energy and serves as a numerical guard; it is therefore kept as a hard max.

\paragraph{Quantitative analysis and biological interpretation.}
Equation \eqref{eq:stim-strong-nd} acts as a spatio-temporal low-pass filter of the mechanostat drive. Regions of sustained overload ($\hat{\psi}>\hat{\psi}_{\mathrm{ref}}$, i.e.\ $\delta>0$) accumulate positive stimulus, while underloaded regions ($\delta<0$) accumulate negative stimulus. Diffusion broadens and smooths these regions over the characteristic length
\(
\ell_S=\sqrt{\tau_S D_S}
\), while the reaction term damps transient fluctuations with a characteristic time
\(
t_S=\tau_S
\).
The time constant $\tau_S$ thus sets the inertia with respect to temporal fluctuations in the drive, while $D_S$ controls spatial spread via gap-junction/vascular pathways. The homogeneous Neumann condition corresponds to zero normal flux of $S$ at $\partial\Omega$.
\rev{We emphasize that the diffusion length $\ell_S = \sqrt{\tau_S D_S}$ is set in practice by a mesh-based prescription $\ell = \alpha h$ (\S\ref{sec:parameter-governance}), yielding $\ell_S$ on the order of $10^{-1}$--$10^{0}$\,mm. This is substantially larger than the osteocyte spacing in the lacuno-canalicular network (${\sim}20$--$50\,\mu$m); accordingly, $D_S$ should be understood as a phenomenological regularization parameter that controls the effective signal reach at the continuum (FE) scale, rather than as a direct representation of LCN-scale diffusion. The saturation width $\kappa_S$ and the lazy-zone threshold $\delta_0$ are likewise phenomenological calibration parameters; we therefore quantify their local sensitivity by a compact full-factorial femur sweep in \S\ref{sec:mechanostat-sensitivity}. Over the tested range, the response is governed primarily by $\kappa_S$, while the default pair $(\delta_0,\kappa_S)=(0.1,0.4)$ does not lie on a sharp optimum or instability boundary.}

\subsection*{Density submodel}\label{subsec:density-nd}
The apparent density $\rho=\rho(x,t)$ evolves under diffusive regularization and bounded formation/resorption kinetics driven by the stimulus $S$. The strong form is supplemented with homogeneous no-flux boundary conditions and an admissible initial state,
\begin{equation}\label{eq:density-bc-ic-nd}
\big(\mathbf{D}_\rho\,\nabla\rho\big)\!\cdot\!\mathbf{n}=0\quad\text{on }\partial\Omega\times(0,T],\qquad
\rho(\cdot,0)=\rho_0\ \ \text{in }\Omega,\ \ \rho_{\min}\le \rho_0\le \rho_{\max}.
\end{equation}

\paragraph{Diffusive regularization.}
In this work, transport is modeled by an isotropic diffusion tensor
\begin{equation}\label{eq:Drho-nd}
\mathbf{D}_\rho(x,t) \;=\; D_\rho\,\mathbf{I},
\end{equation}
with $D_\rho\ge 0$ in units mm$^2$/day.

\paragraph{Bounded formation/resorption kinetics.}
The density evolution is driven by the stimulus $S$ such that positive stimulus drives growth toward $\rho_{\max}$ and negative stimulus drives resorption toward $\rho_{\min}$. We use smooth positive/negative parts (consistent with the global smoothing width $\varepsilon$):
\begin{align}\label{eq:smooth-splits-nd}
S_+(S):=\operatorname{smax}(S,0;\varepsilon),\\
S_-(S):=\operatorname{smax}(-S,0;\varepsilon).
\end{align}
The total evolution equation is:
\begin{equation}\label{eq:density-strong-nd}
\partial_t \rho - \nabla \cdot (\mathbf{D}_\rho \nabla \rho) =
\alpha_{\mathrm{surf}}(\rho)\Big(
k_{\mathrm{form}}\,S_+(S)\Big(1-\frac{\rho}{\rho_{\max}}\Big)
\;+\;k_{\mathrm{res}}\,S_-(S)\Big(1-\frac{\rho}{\rho_{\min}}\Big)
\Big).
\end{equation}
Here $k_{\mathrm{form}},k_{\mathrm{res}}\ge 0$ are formation/resorption gains. The bracketed terms are soft bounds: as $\rho\to\rho_{\max}$ (formation) or $\rho\to\rho_{\min}$ (resorption) the corresponding rate vanishes, so $\rho$ remains naturally attracted to $[\rho_{\min},\rho_{\max}]$ without hard constraints. The optional surface-availability factor $\alpha_{\mathrm{surf}}\in[\alpha_{\min},1]$ scales activity based on a porosity/specific-surface proxy.
\rev{We stress that this closure is not an explicit damage model: in contrast to damage-based constitutive formulations for femur mechanics \revcite{Gizzi2022_DamageBone}, the present framework does not evolve a separate damage internal variable, and the density floor $\rho_{\min}$ merely prevents complete loss of mass and stiffness in the apparent-scale update.}

\paragraph{Surface-availability factor.}
To mimic reduced active remodeling surface in dense bone, we use a smooth factor $\alpha_{\mathrm{surf}}(\rho)\in[\alpha_{\min},1]$ based on a porosity proxy. Let $\rho_t>0$ be the tissue density and define the (dimensionless) porosity-like variable
\begin{align*}
f_{\mathrm{raw}}(\rho) &:= 1-\frac{\rho}{\rho_t},\\
f(\rho) &:= \operatorname{sclamp}(f_{\mathrm{raw}}(\rho), 0, 1; \varepsilon).
\end{align*}
We estimate a specific-surface proxy by blending a Martin-type trabecular polynomial
\[
a_{\mathrm{trab}}(f):=32.3 f-93.9 f^2+134.0 f^3-101.0 f^4+28.8 f^5
\]
with a cortical proxy $a_{\mathrm{cort}}(f):=c_{\mathrm{cort}}\sqrt{f+\varepsilon}$, where
\[
f_{\mathrm{tr}}:=\max\!\Big(1-\frac{\rho_{\mathrm{trab}}^{\max}}{\rho_t},f_{\min}\Big),\qquad
c_{\mathrm{cort}}:=\frac{a_{\mathrm{trab}}(f_{\mathrm{tr}})}{\sqrt{f_{\mathrm{tr}}}},
\]
and $f_{\min}>0$ is a small fixed lower bound. Using the same smooth transition weight $w(\widetilde{\rho})$ as in \eqref{eq:variable-exponent}, define
\[
a_v(\rho):=\operatorname{smax}\!\Big((1-w(\widetilde{\rho}))\,a_{\mathrm{trab}}(f(\rho))+w(\widetilde{\rho})\,a_{\mathrm{cort}}(f(\rho)),\,0;\varepsilon\Big),
\]
and set
\begin{equation}\label{eq:Asurf-nd}
\alpha_{\mathrm{surf}}(\rho):=\alpha_{\min}+(1-\alpha_{\min})\frac{a_v(\rho)}{a_v(\rho)+a_{v,0}},
\end{equation}
with $\alpha_{\min}\in[0,1]$ and $a_{v,0}>0$.

\paragraph{Quantitative analysis and biological interpretation.}
Equation \eqref{eq:density-strong-nd} uses differentiable kinetics and soft bounds to provide a stable density update driven by the nonlocal stimulus field. The diffusion term $\nabla\!\cdot(\mathbf{D}_\rho\nabla\rho)$ introduces a spatial regularization length scale, preventing checkerboard patterns and promoting mesh-independent solutions. Here $\mathbf{D}_\rho$ is isotropic; fabric effects enter through the mechanics and the fabric evolution, and can be extended to anisotropic transport if desired.

\begin{figure}[htbp]
    \centering
    \safeincludegraphics[width=\textwidth]{mechanostat_law.png}
    \caption{Mechanostat and density kinetics. Top row: stimulus mechanostat drive (lazy zone and saturation). Bottom row: surface-availability factor and bounded formation/resorption rates that keep $\rho$ attracted to $[\rho_{\min},\rho_{\max}]$.}
    \label{fig:mechanostat_law}
\end{figure}

\subsection*{Direction (fabric) submodel}\label{subsec:direction-nd}
The local fabric is described by a second-order, symmetric, positive-definite tensor field $\mathbf{A}=\mathbf{A}(x,t)$. To ensure positive definiteness and robustness, we evolve the logarithmic fabric tensor $\mathbf{L} = \ln \mathbf{A}$. The dynamics follow a relaxation--diffusion flow toward a mechanically informed, trace-free logarithmic target $\mathbf{L}_M(\bar{\mathbf{Q}})$ defined below. The strong form on a bounded Lipschitz domain is
\begin{equation}\label{eq:L-strong-nd}
\begin{aligned}
\partial_t \mathbf{L} \;-\; D_A\,\Delta \mathbf{L} \;+\; \frac{a(\bar{\mathbf{Q}})}{\tau_A}\,\mathbf{L}
\;=\; \frac{a(\bar{\mathbf{Q}})}{\tau_A}\,\mathbf{L}_M(\bar{\mathbf{Q}})\,,
\end{aligned}\end{equation}
supplemented with homogeneous no-flux boundary conditions and an admissible initial state,
\begin{equation}\label{eq:L-bc-ic-nd}
\partial_{\mathbf{n}}\mathbf{L}:=\big(\nabla \mathbf{L}\big)\,\mathbf{n}=\mathbf{0}\quad\text{on }\partial\Omega\times(0,T],
\qquad
\mathbf{L}(\cdot,0)=\ln(\mathbf{A}_0)\ \ \text{in }\Omega.
\end{equation}
Here $D_A\ge 0$ is a diffusion coefficient and $\tau_A>0$ a relaxation time constant. The activity factor $a(\bar{\mathbf{Q}})\in[0,1]$ suppresses updates under near-isotropic loading; the Laplacian and normal derivative act componentwise on $\mathbf{L}$.
\rev{At the modelling level, \eqref{eq:L-strong-nd} should be read as an apparent-scale relaxation law. Complementary homogenized constrained-mixture approaches evolve fiber distributions from lower-scale constituent turnover and reorientation \revcite{Vasta2023_ConstrainedMixtureFabric}; here those mechanisms are collapsed into the phenomenological target $\mathbf{L}_M(\bar{\mathbf{Q}})$ and its relaxation rate.}

\paragraph{Target construction from mechanics.}
For each loading case $i$, the mechanics solve provides $\boldsymbol{\sigma}_i$ and $\boldsymbol{\varepsilon}_i$. To extract \emph{directional} information and remove the hydrostatic contribution we form the symmetric deviatoric-stress outer-product
\begin{equation}\label{eq:M-def-nd}
\boldsymbol{\sigma}_{i,\mathrm{dev}}:=\boldsymbol{\sigma}_i-\tfrac13(\mathrm{tr}\,\boldsymbol{\sigma}_i)\mathbf{I},\qquad
\mathbf{Q}_i \;:=\; \mathrm{sym}\!\big(\boldsymbol{\sigma}_{i,\mathrm{dev}}\,\boldsymbol{\sigma}_{i,\mathrm{dev}}^{\!\top}\big),
\end{equation}
and compute a cycle-weighted average
\(
\bar{\mathbf{Q}}:=\sum_i w_i \mathbf{Q}_i\big/\sum_i w_i
\),
where $w_i$ are the same daily weights as in \eqref{eq:stim-psi-nd}. To construct a bounded, trace-free logarithmic target we use the eigen-decomposition of a regularized tensor $\mathbf{Q}:=\mathrm{sym}(\bar{\mathbf{Q}})+\varepsilon_Q\mathbf{I}$, $\mathbf{Q}=\sum_i \lambda_i\mathbf{P}_i$, and define
\begin{equation}\label{eq:M-hat-nd}
\tilde m_i \;:=\; \operatorname{sclamp}((\lambda_i/g)^{\gamma_F}, m_{\min}, m_{\max}; \varepsilon),\qquad
g:=(\lambda_1\lambda_2\lambda_3)^{1/3},\qquad
\mathbf{L}_M(\bar{\mathbf{Q}})\;:=\;\sum_{i=1}^3\Big(\ln(\tilde m_i)-\tfrac13\sum_{j=1}^3\ln(\tilde m_j)\Big)\mathbf{P}_i,
\end{equation}
so that $\mathrm{tr}\,\mathbf{L}_M=0$ (equivalently $\det(\exp(\mathbf{L}_M))=1$) and the induced eigenvalue ratios are bounded by $m_{\min},m_{\max}$. In computation we evaluate $\{\lambda_i,\mathbf{P}_i\}$ using a stabilized Sylvester-projector construction (as above), which is robust under isotropic and transversely isotropic limits of $\mathbf{Q}$.

The activity factor $a(\bar{\mathbf{Q}})$ is computed from the deviatoric magnitude of $\mathbf{Q}$: let $\mathbf{Q}_{\mathrm{dev}}:=\mathbf{Q}-\frac13(\mathrm{tr}\,\mathbf{Q})\mathbf{I}$ and
\[
r:=\frac{\sqrt{\mathbf{Q}_{\mathrm{dev}}:\mathbf{Q}_{\mathrm{dev}}+\varepsilon^2}}{\mathrm{tr}\,\mathbf{Q}+\varepsilon_Q},\qquad
a:=\frac{r}{r+\varepsilon_a}\in[0,1],
\]
with a small dimensionless $\varepsilon_a>0$.

\begin{figure}[htbp]
    \centering
    \safeincludegraphics[width=\textwidth]{anisotropy_law.png}
    \caption{Fabric-guided anisotropy. (a) Directional stiffness for representative fabric states. (b,d) Directional diffusivity illustrations correspond to an optional anisotropic-transport extension; in the present model we use isotropic density diffusion $\mathbf{D}_\rho=D_\rho\mathbf{I}$ in \eqref{eq:Drho-nd}. (c,e,f) Stiffness scaling and eigenvalue-space visualizations.}
    \label{fig:anisotropy_law}
\end{figure}

\paragraph{Quantitative analysis and biological interpretation.}
Figure \ref{fig:anisotropy_law} visualizes how the log-fabric state biases directional stiffness. The relaxation term in \eqref{eq:L-strong-nd} drives $\mathbf{L}$ toward the mechanically informed target $\mathbf{L}_M(\bar{\mathbf{Q}})$ while diffusion regularizes spatial oscillations; the log-Euclidean representation ensures $\mathbf{A}=\exp(\mathbf{L})\in\mathbb{S}^{3}_{++}$ by construction. The log-space normalization in \eqref{eq:M-hat-nd} enforces $\mathrm{tr}\,\mathbf{L}_M=0$ (unit determinant for $\exp(\mathbf{L}_M)$), so the target acts on anisotropy (eigenvalue ratios) rather than isotropic scaling. In the constitutive evaluation, isotropic scaling is removed by working with eigenvalue ratios (deviatoric part of $\mathbf{L}$) and bounding them by the same $m_{\min},m_{\max}$ used in \eqref{eq:M-hat-nd}.
